\section*{Chapter \ref{ch:hf:retail-wholesaling} --- Retail Wholesaling}

\begin{solution}{exo:hf:retail-wholesaling:1}
The half-spread is one cent; 30\% of it is 0.3 cent: the buy fills at \$50.017. The effective spread is $2\times(50.017-50.010)=1.4$ cents, 70\% of
the quoted two.
\end{solution}

\begin{solution}{exo:hf:retail-wholesaling:2}
$1.25(1-\pi)=0.1+0.1+0.2$: $\pi=0.68$.
\end{solution}

\begin{solution}{exo:hf:retail-wholesaling:3}
E/Q $=1-\pi=0.70$: both spreads scale with the half-spread.
\end{solution}

\begin{solution}{exo:hf:retail-wholesaling:4}
Without herding, buys and sells alternate at random and the inventory wanders slowly; with herding, runs of buys or sells push it past the limit, and
the excess must be hedged: 12.9\% of the volume against 0.9\%.
\end{solution}

\begin{solution}{exo:hf:retail-wholesaling:5}
Institutional orders carry 1.5 cents of information a share in the model, more than the half-spread: filling them at the quote loses 0.57 cent a
share before any improvement.
\end{solution}

\begin{solution}{exo:hf:retail-wholesaling:6}
Differences in the orders' types, sizes and timing across brokers, the wholesalers' pricing agreements with each broker (such as improvement
commitments), and differences in the brokers' own routing.
\end{solution}

\begin{solution}{exo:hf:retail-wholesaling:7}
The mean net is highest at a limit of 10,000 shares (0.45 cent); the mean less the standard deviation is highest at 5,000 (0.39 less 0.21).
\end{solution}

\begin{solution}{exo:hf:retail-wholesaling:8}
E/Q measures improvement against the quote at arrival, not the price the customer could have obtained elsewhere; the same order can receive different
prices through different brokers, and the payment for flow is not in the statistic.
\end{solution}

\begin{solution}{pb:hf:retail-wholesaling:1}
\begin{enumerate}
\item $h(1-\pi)-a-p-c_H+I$: capture, mark-out, payment for flow, hedging, inventory P\&L.
\item See \cref{def:hf:retail-wholesaling:segmentation}.
\item Retail: 20,000 orders, median 50 shares, 0.1 cent of information, herding; institutional: median 2,000 shares, 1.5 cents.
\item The improvement to the trader, the payment to the broker and the wholesaler's margin.
\item Retail: 0.85, 0.475 and $-0.025$ cent a share; institutional: $-0.57$, $-0.94$, $-1.44$.
\item At 68\% of the half-spread.
\item See \cref{def:hf:retail-wholesaling:rate}; 87\%.
\item Quoted 2.5 cents, effective 1.75, E/Q 0.70, improvement 0.375 cent a share.
\item Hedging costs against inventory risk: 0.20 net with 0.06 s.d.\ at 1,000 shares, 0.45 with 0.38 at 10,000, dominated by inventory at 50,000
(s.d.\ 1.32).
\item It makes the flow one-sided in waves: 12.9\% of volume hedged against 0.9\% without it.
\item The SEC withdrew the order competition and best-execution proposals in June 2025; Rule 605's amendments now apply from August 2026; the NYSE's
retail programme was approved in 2012 with sub-penny improvement.
\item See \cref{def:hf:retail-wholesaling:rlp}.
\item At 30\% improvement and 0.1 cent paid for flow, the wholesaler keeps 0.475 cent a share before inventory P\&L (0.44 after), and the trader
receives 0.375 cent a share against a 2.5-cent quoted spread (E/Q 0.70).
\item All three: the model splits 0.95 cent a share into 0.375, 0.1 and 0.475.
\item Wholesalers priced the same orders differently by broker; disclosures should report prices relative to comparable executions, by broker.
\item Many retail orders would have been exposed to auctions where others could compete.
\item Equity retail internalisation.
\item From the flow's mark-outs by broker and the order sizes and timing.
\item Its improvement and fill statistics on the broker's own orders, against other wholesalers'.
\item The right to trade against flow that knows little, and the obligation to price it well.
\end{enumerate}
\end{solution}

\begin{solution}{iq:hf:retail-wholesaling:1}
The wholesaler knows the flow is retail and nearly uninformed; exchange quotes must protect against everyone, including informed traders, so they are
wider.
\iqlookfor{segmentation.}
\end{solution}

\begin{solution}{iq:hf:retail-wholesaling:2}
Mark-outs of its orders at several horizons, by order size, type and time, against a benchmark flow.
\iqlookfor{mark-outs by segment.}
\end{solution}

\begin{solution}{iq:hf:retail-wholesaling:3}
Quotes from a consolidated feed, per-broker terms as a lookup table, the inventory as shared state; a function of those three computed on each order,
with limits checked in the same path.
\iqlookfor{precomputed terms, shared inventory.}
\end{solution}

\begin{solution}{iq:hf:retail-wholesaling:4}
Per-name inventory and loss limits; hedging beyond them on exchanges; wider or no improvement in the name, and routing to exchanges if the
inventory cannot be hedged.
\iqlookfor{limits and graceful withdrawal.}
\end{solution}

\begin{solution}{iq:hf:retail-wholesaling:5}
From the net per share of that broker's flow after improvement, hedging and inventory: pay only what leaves a margin, and compare with the
improvement the broker's customers receive.
\iqlookfor{per-broker economics.}
\end{solution}

\begin{solution}{iq:hf:retail-wholesaling:6}
For an AR(1) sign process the variance of the sum grows like $N(1+\rho)/(1-\rho)$ for large $N$: persistence multiplies the imbalance's variance.
\iqlookfor{long-run variance of an AR(1).}
\end{solution}
